% ============================================================
% CHAPITRE 19 — Typographie et lettrage dans Illustrator
% ============================================================

\chapter{Typographie et lettrage dans Illustrator}

\noindent
\textit{Dans les chapitres précédents, nous avons appris à dessiner, à manipuler des formes, à appliquer des couleurs, des dégradés, des contours. Nous savons désormais créer des illustrations vectorielles précises. Mais une illustration sans texte est souvent incomplète : c'est le texte qui porte le message, identifie la marque, structure la composition. Il faut donc apprendre à maîtriser la typographie dans Illustrator.}

\vspace{1em}

\noindent
La \terme{typographie} est l'un des piliers du design graphique. Nous l'avons étudiée en théorie dans le chapitre 4 : anatomie des caractères, familles de polices, lisibilité, interlignage, crénage, hiérarchie. Nous allons maintenant la mettre en pratique dans Illustrator, qui est l'un des meilleurs outils pour la typographie vectorielle. Illustrator permet de créer, modifier, composer, déformer du texte avec une précision et une liberté remarquables.

\vspace{1em}

\noindent
Mais Illustrator ne se limite pas à la typographie standard : il permet aussi le \terme{lettrage} (ou \emph{lettering}), c'est-à-dire la création de lettres personnalisées, dessinées à la main, uniques. Le lettrage est utilisé pour les logos, les titres, les affiches, les identités visuelles. Il permet de créer des caractères sur mesure, adaptés à un projet, à une ambiance, à une marque. C'est l'une des compétences les plus recherchées chez un designer graphique.

\vspace{1em}

\noindent
Ce chapitre a pour objectif de te faire maîtriser la typographie et le lettrage dans Illustrator. Nous commencerons par apprendre à créer et modifier du texte. Nous étudierons ensuite les propriétés typographiques, le crénage et l'interlignage, les compositions typographiques, le texte sur tracé, la vectorisation, la personnalisation des lettres, le lettrage simple, la conception d'un titre publicitaire, et la préservation des versions modifiables. À la fin de ce chapitre, tu sauras créer des compositions typographiques professionnelles.

% ============================================================
\section{Créer et modifier du texte}
% ============================================================

\subsection{Pourquoi créer du texte dans Illustrator ?}

Illustrator n'est pas le logiciel idéal pour la typographie éditoriale (InDesign est plus adapté pour les documents longs). Mais Illustrator est très utilisé pour :

\begin{itemize}
    \item \terme{Créer} des logos typographiques.
    \item \terme{Composer} des titres, des slogans, des accroches.
    \item \terme{Intégrer} du texte dans une illustration.
    \item \terme{Dessiner} du lettrage personnalisé.
    \item \terme{Préparer} des affiches, des couvertures, des packagings.
\end{itemize}

Illustrator propose des outils de texte complets, qui permettent de créer, modifier, styliser, déformer du texte avec une grande précision.

\subsection{Les outils de texte}

\subsubsection{L'outil Texte}

L'\terme{outil Texte} (\texttt{T}) permet de créer du texte. Il propose plusieurs modes.

\subsubsection{L'outil Texte vertical}

L'\terme{outil Texte vertical} permet de créer du texte vertical (caractères empilés verticalement). Il est utilisé pour des usages spécifiques (affiches japonaises, titres verticaux, colonnes).

\subsubsection{L'outil Texte sur tracé}

L'\terme{outil Texte sur tracé} permet de créer du texte qui suit un tracé (courbe, forme). Il est utilisé pour les logos circulaires, les badges, les textes curvilignes.

\subsubsection{L'outil Texte de paragraphe}

L'\terme{outil Texte de paragraphe} permet de créer un bloc de texte contenu dans une zone. Il est utilisé pour les paragraphes, les descriptions.

\subsubsection{L'outil Texte curviligne}

L'\terme{outil Texte curviligne} permet de créer du texte qui suit une courbe.

\subsection{Les types de texte}

\subsubsection{Le texte ponctuel}

Le \terme{texte ponctuel} est un texte créé par un simple clic sur le plan de travail. Il n'a pas de zone délimitée : il s'étend sur une ligne (ou plusieurs lignes si on appuie sur \texttt{Entrée}). Il est utilisé pour :

\begin{itemize}
    \item Les titres.
    \item Les slogans.
    \item Les mots isolés.
    \item Les étiquettes.
\end{itemize}

\subsubsection{Le texte de paragraphe}

Le \terme{texte de paragraphe} est un texte créé en cliquant-glissant pour définir une zone rectangulaire. Il s'étend dans cette zone, avec des retours à la ligne automatiques. Il est utilisé pour :

\begin{itemize}
    \item Les paragraphes.
    \item Les descriptions.
    \item Les mentions.
    \item Les blocs de contenu.
\end{itemize}

\subsubsection{Le texte curviligne}

Le \terme{texte curviligne} est un texte créé sur un tracé. Il suit la courbe du tracé. Il est utilisé pour :

\begin{itemize}
    \item Les logos circulaires.
    \item Les badges.
    \item Les tampons.
    \item Les compositions créatives.
\end{itemize}

\subsection{Créer un texte ponctuel}

\subsubsection{Étapes}

\begin{enumerate}
    \item Sélectionner l'outil Texte (\texttt{T}).
    \item Cliquer sur le plan de travail.
    \item Saisir le texte.
    \item Valider avec \texttt{Échap} ou en cliquant ailleurs.
\end{enumerate}

\subsubsection{Le calque de texte}

Un calque de texte est automatiquement créé. Il est reconnaissable à l'icône \textbf{T} dans le panneau Calques. Le texte est modifiable tant qu'il n'est pas vectorisé.

\subsection{Créer un texte de paragraphe}

\subsubsection{Étapes}

\begin{enumerate}
    \item Sélectionner l'outil Texte (\texttt{T}).
    \item Cliquer-glisser pour définir une zone rectangulaire.
    \item Saisir le texte.
    \item Valider.
\end{enumerate}

\subsubsection{Le cadre de texte}

Le texte de paragraphe est contenu dans un cadre. On peut :

\begin{itemize}
    \item \terme{Redimensionner} le cadre en tirant sur les poignées.
    \item \terme{Déplacer} le cadre en cliquant-glissant à l'intérieur.
    \item \terme{Dépasser} le cadre : si le texte est trop long, un signe \textbf{+} apparaît en bas à droite.
\end{itemize}

\subsection{Créer un texte curviligne}

\subsubsection{Étapes}

\begin{enumerate}
    \item Créer un tracé (avec la Plume, ou une forme).
    \item Sélectionner l'outil Texte sur tracé.
    \item Cliquer sur le tracé.
    \item Saisir le texte.
    \item Valider.
\end{enumerate}

\subsubsection{Le tracé de texte}

Le texte suit le tracé. On peut :

\begin{itemize}
    \item \terme{Déplacer} le texte le long du tracé.
    \item \terme{Retourner} le texte (de l'autre côté du tracé).
    \item \terme{Modifier} le tracé pour changer la courbe.
\end{itemize}

\subsection{Modifier un texte}

\subsubsection{Modifier le contenu}

\begin{enumerate}
    \item Sélectionner l'outil Texte (\texttt{T}).
    \item Cliquer sur le texte existant.
    \item Modifier le contenu.
    \item Valider.
\end{enumerate}

\subsubsection{Modifier les propriétés}

\begin{itemize}
    \item \terme{Police} : changer la police.
    \item \terme{Taille} : changer la taille.
    \item \terme{Couleur} : changer la couleur.
    \item \terme{Graisse} : changer la graisse (normal, gras, etc.).
    \item \terme{Style} : changer le style (italique, souligné).
    \item \terme{Alignement} : changer l'alignement (gauche, centre, droite, justifié).
\end{itemize}

\subsection{Vectoriser un texte}

\subsubsection{Principe}

\terme{Vectoriser} un texte consiste à le transformer en formes vectorielles. Le texte n'est plus modifiable, mais il peut être déformé, stylisé, transformé comme une forme.

\subsubsection{Utilisation}

\texttt{Texte > Vectoriser} (ou \texttt{Ctrl + Maj + O}).

\subsubsection{Précautions}

\begin{itemize}
    \item \terme{Conserver} une copie du calque de texte avant de vectoriser.
    \item \terme{Vectoriser} uniquement si nécessaire.
    \item \terme{Ne pas vectoriser} un texte destiné à être modifié.
\end{itemize}

\subsection{Les bonnes pratiques}

\begin{itemize}
    \item \terme{Utiliser} l'outil Texte pour créer du texte.
    \item \terme{Choisir} le type de texte adapté (ponctuel, paragraphe, curviligne).
    \item \terme{Nommer} le calque de texte.
    \item \terme{Conserver} le texte modifiable (ne pas vectoriser sans nécessité).
    \item \terme{Vérifier} la lisibilité (taille, contraste).
    \item \terme{Utiliser} les polices adaptées.
    \item \terme{Enregistrer} régulièrement.
\end{itemize}

\subsection*{Conclusion de la section}

Créer du texte dans Illustrator est une opération courante. On utilise l'outil Texte (\texttt{T}) pour créer du texte ponctuel, du texte de paragraphe, ou du texte curviligne. Le texte est un calque spécifique, modifiable tant qu'il n'est pas vectorisé. On peut modifier le contenu, les propriétés, l'alignement. Il est essentiel de nommer les calques, de conserver le texte modifiable, de vérifier la lisibilité.

Maintenant que nous savons créer du texte, nous allons, dans la section suivante, apprendre à \terme{comprendre les propriétés typographiques}.

% ============================================================
\section{Comprendre les propriétés typographiques}
% ============================================================

\subsection{Les panneaux de typographie}

\subsubsection{Le panneau Caractère}

Le panneau \terme{Caractère} (\texttt{Fenêtre > Texte > Caractère}) permet de modifier les propriétés typographiques du texte sélectionné :

\begin{itemize}
    \item \terme{Police} : choisir la police.
    \item \terme{Style} : normal, gras, italique, etc.
    \item \terme{Taille} : taille du texte.
    \item \terme{Interlignage} : espacement entre les lignes.
    \item \terme{Crénage} : espacement entre les caractères.
    \item \terme{Approche} : espacement global entre les caractères.
    \item \terme{Décalage vertical} : décalage de la ligne de base.
    \item \terme{Décalage horizontal} : décalage horizontal.
    \item \terme{Rotation} : rotation des caractères.
    \item \terme{Couleur} : couleur du texte.
    \item \terme{Majuscules/Minuscules} : petites majuscules, exposant, indice.
    \item \terme{Soulignement, barré} : styles de décoration.
    \item \terme{Langue} : langue du texte (pour la césure et la correction).
\end{itemize}

\subsubsection{Le panneau Paragraphe}

Le panneau \terme{Paragraphe} (\texttt{Fenêtre > Texte > Paragraphe}) permet de modifier les propriétés de mise en page du texte :

\begin{itemize}
    \item \terme{Alignement} : gauche, centre, droite, justifié.
    \item \terme{Retrait} : retrait de première ligne, retrait à gauche, retrait à droite.
    \item \terme{Espacement} : espacement avant et après les paragraphes.
    \item \terme{Césure} : césure automatique.
    \item \terme{Ligne de base} : décalage de la ligne de base.
    \item \terme{Composition} : options de composition.
\end{itemize}

\subsection{La barre de contrôle}

Lorsque l'outil Texte est sélectionné, la barre de contrôle propose :

\begin{itemize}
    \item \terme{Police} : choisir la police.
    \item \terme{Style} : normal, gras, italique.
    \item \terme{Taille} : taille du texte.
    \item \terme{Alignement} : gauche, centre, droite, justifié.
    \item \terme{Couleur} : couleur du texte.
    \item \terme{Déformation} : déformer le texte.
    \item \terme{Panoramique} : créer un texte.
    \item \terme{Paragraphe} : créer un texte de paragraphe.
\end{itemize}

\subsection{La couleur du texte}

\subsubsection{Changer la couleur}

\begin{enumerate}
    \item Sélectionner le texte.
    \item Choisir une couleur dans la case de remplissage (barre d'outils).
    \item Ou utiliser le panneau \emph{Couleur} ou \emph{Nuancier}.
\end{enumerate}

\subsubsection{Appliquer une couleur différente à une partie du texte}

\begin{enumerate}
    \item Sélectionner une partie du texte avec l'outil Texte.
    \item Changer la couleur.
    \item Valider.
\end{enumerate}

\subsection{L'interlignage, le crénage, l'approche}

\subsubsection{L'interlignage}

L'\terme{interlignage} est l'espace vertical entre les lignes. Il se mesure en points ou en pourcentage de la taille du texte.

\begin{itemize}
    \item \terme{Interlignage automatique} : calculé par Illustrator.
    \item \terme{Interlignage manuel} : défini par l'utilisateur.
    \item \terme{Valeurs courantes} : 120\% à 150\% de la taille du texte.
\end{itemize}

\subsubsection{Le crénage}

Le \terme{crénage} est l'espacement entre deux caractères spécifiques. Il se mesure en millièmes d'em.

\begin{itemize}
    \item \terme{Crénage métrique} : selon les métriques de la police.
    \item \terme{Crénage optique} : ajustement automatique.
    \item \terme{Crénage manuel} : ajustement manuel.
\end{itemize}

\subsubsection{L'approche}

L'\terme{approche} est l'espacement global entre tous les caractères. Elle se mesure en millièmes d'em.

\begin{itemize}
    \item \terme{Approche négative} : les caractères sont rapprochés.
    \item \terme{Approche positive} : les caractères sont espacés.
\end{itemize}

\subsection{Les styles de texte}

\subsubsection{Les styles de paragraphe}

Les \terme{styles de paragraphe} permettent d'enregistrer des propriétés de texte pour les réutiliser. On peut créer un style pour :

\begin{itemize}
    \item Les titres.
    \item Les sous-titres.
    \item Le texte courant.
    \item Les légendes.
\end{itemize}

\subsubsection{Les styles de caractère}

Les \terme{styles de caractère} permettent d'enregistrer des propriétés de caractère (police, taille, couleur) pour les réutiliser.

\subsubsection{Création d'un style}

\begin{enumerate}
    \item Sélectionner le texte avec les propriétés souhaitées.
    \item Ouvrir le panneau \emph{Styles de paragraphe} ou \emph{Styles de caractère}.
    \item Cliquer sur \emph{Créer un style}.
    \item Nommer le style.
    \item Valider.
\end{enumerate}

\subsection{Les bonnes pratiques}

\begin{itemize}
    \item \terme{Choisir} des polices lisibles.
    \item \terme{Limiter} le nombre de polices (1 à 2).
    \item \terme{Utiliser} des tailles cohérentes (échelle typographique).
    \item \terme{Ajuster} l'interlignage pour la lisibilité.
    \item \terme{Utiliser} des styles pour la cohérence.
    \item \terme{Vérifier} la lisibilité sur différents fonds.
    \item \terme{Ne pas abuser} des effets.
\end{itemize}

\subsection*{Conclusion de la section}

Les \terme{propriétés typographiques} permettent de contrôler l'apparence du texte : police, taille, style, interlignage, crénage, approche, couleur, alignement. On les modifie via la barre de contrôle, les panneaux Caractère et Paragraphe. On peut enregistrer des styles pour la cohérence. Il est essentiel de choisir des polices lisibles, de limiter le nombre de polices, d'ajuster l'interlignage, d'utiliser des styles, de vérifier la lisibilité.

Maintenant que nous comprenons les propriétés typographiques, nous allons, dans la section suivante, apprendre à \terme{régler le crénage et l'interlignage}.

% ============================================================
\section{Régler le crénage et l'interlignage}
% ============================================================

\subsection{L'importance du crénage et de l'interlignage}

Le \terme{crénage} et l'\terme{interlignage} sont deux paramètres essentiels de la typographie. Ils déterminent :

\begin{itemize}
    \item \terme{La lisibilité} du texte.
    \item \terme{L'esthétique} de la composition.
    \item \terme{La qualité} du rendu professionnel.
    \item \terme{La personnalité} du texte.
\end{itemize}

Un texte mal créné ou mal interligné est désagréable à lire, même avec une belle police. Un texte bien créné et bien interligné est agréable, lisible, professionnel.

\subsection{Le crénage}

\subsubsection{Définition}

Le \terme{crénage} (ou \emph{kerning}) est l'ajustement de l'espace entre deux caractères spécifiques. Il permet de corriger les espacements irréguliers.

\subsubsection{Pourquoi ajuster le crénage ?}

Certaines paires de lettres créent des espacements visuellement irréguliers :

\begin{itemize}
    \item \terme{AV} : espace trop large.
    \item \terme{To} : espace trop large.
    \item \terme{WA} : espace trop large.
    \item \terme{fi} : espace trop serré.
\end{itemize}

L'ajustement du crénage permet d'obtenir un espacement régulier.

\subsubsection{Les méthodes de crénage}

\begin{itemize}
    \item \terme{Crénage métrique} : utilise les métriques de la police (par défaut).
    \item \terme{Crénage optique} : ajustement automatique par Illustrator.
    \item \terme{Crénage manuel} : ajustement manuel avec les flèches ou les raccourcis.
\end{itemize}

\subsubsection{Comment ajuster le crénage ?}

\begin{enumerate}
    \item Placer le curseur entre les deux caractères.
    \item Ouvrir le panneau \emph{Caractère}.
    \item Ajuster la valeur de crénage (en millièmes d'em).
    \item Ou utiliser \texttt{Alt + flèches} (gauche/droite) pour ajuster manuellement.
\end{enumerate}

\subsection{L'interlignage}

\subsubsection{Définition}

L'\terme{interlignage} (ou \emph{leading}) est l'espace vertical entre les lignes. Il se mesure de ligne de base à ligne de base.

\subsubsection{Pourquoi ajuster l'interlignage ?}

L'interlignage détermine :

\begin{itemize}
    \item \terme{La lisibilité} : un interlignage trop serré fatigue, un interlignage trop large disperse.
    \item \terme{L'esthétique} : un interlignage bien ajusté donne un texte équilibré.
    \item \terme{La densité} : un interlignage serré donne un texte dense, un interlignage large donne un texte aéré.
\end{itemize}

\subsubsection{Les valeurs recommandées}

\begin{itemize}
    \item \terme{Texte courant} : 120\% à 150\% de la taille du texte.
    \item \terme{Titres} : 100\% à 120\%.
    \item \terme{Petits textes} : 130\% à 160\%.
\end{itemize}

\subsubsection{Comment ajuster l'interlignage ?}

\begin{enumerate}
    \item Sélectionner le texte.
    \item Ouvrir le panneau \emph{Caractère}.
    \item Saisir une valeur d'interlignage (en points ou en pourcentage).
    \item Ou utiliser \texttt{Alt + flèches} (haut/bas) pour ajuster manuellement.
\end{enumerate}

\subsection{Le crénage et l'interlignage en pratique}

\subsubsection{Pour les titres}

\begin{itemize}
    \item \terme{Crénage} : ajusté manuellement (les titres sont grands, les défauts sont visibles).
    \item \terme{Interlignage} : serré (100\% à 120\%) pour un effet compact.
\end{itemize}

\subsubsection{Pour le texte courant}

\begin{itemize}
    \item \terme{Crénage} : automatique (métrique ou optique).
    \item \terme{Interlignage} : 120\% à 150\%.
\end{itemize}

\subsubsection{Pour les logos}

\begin{itemize}
    \item \terme{Crénage} : ajusté caractère par caractère (précision extrême).
    \item \terme{Interlignage} : ajusté visuellement.
\end{itemize}

\subsection{Les bonnes pratiques}

\begin{itemize}
    \item \terme{Ajuster} le crénage pour les titres et les logos.
    \item \terme{Ajuster} l'interlignage pour la lisibilité.
    \item \terme{Utiliser} le crénage optique pour le texte courant.
    \item \terme{Vérifier} le rendu à taille réelle.
    \item \terme{Ne pas trop serrer} l'interlignage (perte de lisibilité).
    \item \terme{Ne pas trop espacer} le crénage (perte de cohérence).
\end{itemize}

\subsection*{Conclusion de la section}

Le \terme{crénage} et l'\terme{interlignage} sont deux paramètres essentiels de la typographie. Le crénage ajuste l'espace entre deux caractères spécifiques ; l'interlignage ajuste l'espace entre les lignes. On les ajuste via le panneau Caractère ou les raccourcis (\texttt{Alt + flèches}). Il est essentiel d'ajuster le crénage pour les titres et les logos, l'interlignage pour la lisibilité, de vérifier le rendu à taille réelle.

Maintenant que nous savons régler le crénage et l'interlignage, nous allons, dans la section suivante, apprendre à \terme{créer des compositions typographiques}.

% ============================================================
\section{Créer des compositions typographiques}
% ============================================================

\subsection{Qu'est-ce qu'une composition typographique ?}

\begin{quote}
\textbf{Une composition typographique est un agencement réfléchi de texte (titres, sous-titres, paragraphes, mots isolés), qui combine police, taille, graisse, couleur, position, espacement, pour transmettre un message et créer un effet visuel.}
\end{quote}

Une composition typographique ne se limite pas à écrire du texte : elle organise le texte pour créer du sens, de l'émotion, de l'impact.

\subsection{Les principes de la composition typographique}

\subsubsection{La hiérarchie}

La \terme{hiérarchie} organise les niveaux de texte selon leur importance :

\begin{itemize}
    \item \terme{Titre principal} : le plus grand, le plus visible.
    \item \terme{Sous-titre} : secondaire.
    \item \terme{Texte courant} : plus petit.
    \item \terme{Légendes} : encore plus petit.
    \item \terme{Mentions} : le plus petit.
\end{itemize}

\subsubsection{L'alignement}

L'\terme{alignement} organise le texte :

\begin{itemize}
    \item \terme{Aligné à gauche} : le plus lisible.
    \item \terme{Centré} : formel, pour les titres.
    \item \terme{Aligné à droite} : rare, pour des effets.
    \item \terme{Justifié} : élégant, pour les blocs.
\end{itemize}

\subsubsection{Le contraste}

Le \terme{contraste} distingue les éléments :

\begin{itemize}
    \item \terme{Contraste de taille} : grand/petit.
    \item \terme{Contraste de graisse} : gras/fin.
    \item \terme{Contraste de couleur} : vif/neutre.
    \item \terme{Contraste de police} : serif/sans-serif.
\end{itemize}

\subsubsection{L'espace}

L'\terme{espace} aère la composition :

\begin{itemize}
    \item \terme{Espace entre les lignes} : interlignage.
    \item \terme{Espace entre les paragraphes} : respiration.
    \item \terme{Espace autour du texte} : marges.
\end{itemize}

\subsection{Les types de compositions typographiques}

\subsubsection{La composition classique}

La \terme{composition classique} utilise des polices traditionnelles, des alignements symétriques, des hiérarchies claires. Elle est utilisée pour les documents formels, les invitations, les livres.

\subsubsection{La composition moderne}

La \terme{composition moderne} utilise des polices sans-serif, des alignements asymétriques, des contrastes forts. Elle est utilisée pour les affiches, les publicités, les sites web.

\subsubsection{La composition expressive}

La \terme{composition expressive} utilise des polices originales, des déformations, des superpositions, des effets. Elle est utilisée pour les affiches créatives, les couvertures, les visuels d'art.

\subsubsection{La composition minimaliste}

La \terme{composition minimaliste} utilise peu d'éléments, des espaces larges, des polices simples. Elle est utilisée pour les marques de luxe, les identités épurées, les sites modernes.

\subsection{Construire une composition typographique}

\subsubsection{Étape 1 : Définir le message}

Quel est le message principal ? Que doit retenir le lecteur ?

\subsubsection{Étape 2 : Définir la hiérarchie}

Quels sont les niveaux de texte ? Quel est l'ordre d'importance ?

\subsubsection{Étape 3 : Choisir les polices}

Quelles polices utiliser ? 1 à 2 polices maximum.

\subsubsection{Étape 4 : Choisir les tailles}

Quelles tailles pour chaque niveau ? Échelle typographique.

\subsubsection{Étape 5 : Choisir les couleurs}

Quelles couleurs pour chaque niveau ? Cohérence et contraste.

\subsubsection{Étape 6 : Composer}

Placer les éléments selon la hiérarchie, l'alignement, l'espace.

\subsubsection{Étape 7 : Vérifier}

Vérifier la lisibilité, la cohérence, l'impact.

\subsection{Les bonnes pratiques}

\begin{itemize}
    \item \terme{Limiter} le nombre de polices (1 à 2).
    \item \terme{Définir} une hiérarchie claire.
    \item \terme{Utiliser} des tailles cohérentes.
    \item \terme{Respecter} l'espace négatif.
    \item \terme{Créer} du contraste.
    \item \terme{Vérifier} la lisibilité.
    \item \terme{Tester} sur différents supports.
    \item \terme{Faire} une pause avant de finaliser.
\end{itemize}

\subsection*{Conclusion de la section}

Une \terme{composition typographique} est un agencement réfléchi de texte. Elle repose sur la hiérarchie, l'alignement, le contraste, l'espace. Elle peut être classique, moderne, expressive, minimaliste. La construction suit plusieurs étapes : définir le message, la hiérarchie, les polices, les tailles, les couleurs, composer, vérifier. Il est essentiel de limiter les polices, de définir une hiérarchie, de respecter l'espace, de vérifier la lisibilité.

Maintenant que nous savons créer des compositions typographiques, nous allons, dans la section suivante, apprendre à \terme{utiliser le texte sur un tracé}.

% ============================================================
\section{Utiliser le texte sur un tracé}
% ============================================================

\subsection{Qu'est-ce que le texte sur tracé ?}

\begin{quote}
\textbf{Le texte sur tracé (ou \emph{text on a path}) est un texte qui suit la forme d'un tracé (courbe, forme, ligne). Il permet de créer des textes curvilignes, des logos circulaires, des badges, des compositions créatives.}
\end{quote}

Le texte sur tracé est utilisé pour :

\begin{itemize}
    \item \terme{Logos} circulaires.
    \item \terme{Badges} et tampons.
    \item \terme{Couvertures} de magazines.
    \item \terme{Affiches} créatives.
    \item \terme{Illustrations} avec texte.
\end{itemize}

\subsection{Créer un texte sur tracé}

\subsubsection{Avec l'outil Texte sur tracé}

\begin{enumerate}
    \item Créer un tracé (avec la Plume, ou une forme).
    \item Sélectionner l'outil \emph{Texte sur tracé} (dans le groupe de l'outil Texte).
    \item Cliquer sur le tracé.
    \item Saisir le texte.
    \item Valider.
\end{enumerate}

\subsubsection{Avec l'outil Texte standard}

\begin{enumerate}
    \item Créer un tracé.
    \item Sélectionner l'outil Texte (\texttt{T}).
    \item Approcher le curseur du tracé (le curseur change).
    \item Cliquer sur le tracé.
    \item Saisir le texte.
\end{enumerate}

\subsection{Le tracé de texte}

\subsubsection{Le point de départ}

Le texte commence au point où on a cliqué sur le tracé. On peut déplacer le texte le long du tracé.

\subsubsection{Déplacer le texte}

\begin{enumerate}
    \item Sélectionner l'outil \emph{Sélection directe} (\texttt{A}).
    \item Placer le curseur sur le texte (le curseur change).
    \item Glisser le texte le long du tracé.
\end{enumerate}

\subsubsection{Retourner le texte}

\begin{enumerate}
    \item Sélectionner l'outil \emph{Sélection directe}.
    \item Glisser le texte de l'autre côté du tracé.
    \item Le texte se retourne.
\end{enumerate}

\subsection{Les options du texte sur tracé}

\subsubsection{Les options d'effet}

\begin{itemize}
    \item \terme{Aligner sur le tracé} : le texte est aligné sur le tracé.
    \item \terme{Décalage de ligne de base} : le texte est décalé par rapport au tracé.
    \item \terme{Retourner} : le texte est retourné.
\end{itemize}

\subsubsection{Les options de paragraphe}

\begin{itemize}
    \item \terme{Alignement} : le texte est aligné sur le tracé (gauche, centre, droite).
    \item \terme{Espacement} : l'espacement entre les caractères et les lignes.
\end{itemize}

\subsection{Modifier le tracé}

\begin{enumerate}
    \item Sélectionner l'outil \emph{Sélection directe}.
    \item Modifier le tracé (déplacer les points, ajuster les poignées).
    \item Le texte suit le tracé modifié.
\end{enumerate}

\subsection{Les utilisations}

Le texte sur tracé est utilisé pour :

\begin{itemize}
    \item \terme{Logos} circulaires (autour d'un cercle).
    \item \terme{Badges} (texte courbe).
    \item \terme{Couvertures} (titres courbes).
    \item \terme{Affiches} (compositions créatives).
    \item \terme{Illustrations} (texte intégré).
\end{itemize}

\subsection{Les bonnes pratiques}

\begin{itemize}
    \item \terme{Créer} un tracé adapté (cercle, courbe).
    \item \terme{Placer} le texte au bon endroit.
    \item \terme{Ajuster} le décalage de ligne de base.
    \item \terme{Vérifier} la lisibilité.
    \item \terme{Modifier} le tracé si nécessaire.
    \item \terme{Éviter} les courbes trop serrées (lisibilité).
\end{itemize}

\subsection*{Conclusion de la section}

Le \terme{texte sur tracé} est un texte qui suit la forme d'un tracé. On le crée avec l'outil Texte sur tracé, ou l'outil Texte standard. On peut déplacer le texte le long du tracé, le retourner, ajuster le décalage. Il est utilisé pour les logos circulaires, les badges, les affiches. Il est essentiel de créer un tracé adapté, de placer le texte au bon endroit, de vérifier la lisibilité.

Maintenant que nous savons utiliser le texte sur tracé, nous allons, dans la section suivante, apprendre à \terme{vectoriser du texte lorsque nécessaire}.

% ============================================================
\section{Vectoriser du texte lorsque nécessaire}
% ============================================================

\subsection{Qu'est-ce que la vectorisation de texte ?}

\begin{quote}
\textbf{La vectorisation de texte consiste à transformer un texte en formes vectorielles. Le texte n'est plus modifiable comme du texte, mais il devient un ensemble de tracés, modifiables comme des formes.}
\end{quote}

La vectorisation est une opération importante, qui a des avantages et des inconvénients.

\subsection{Quand vectoriser un texte ?}

\subsubsection{Les cas où la vectorisation est nécessaire}

\begin{itemize}
    \item \terme{Envoyer} un fichier à un imprimeur (pour éviter les problèmes de polices).
    \item \terme{Créer} un logo (les lettres deviennent des formes).
    \item \terme{Personnaliser} un texte (modifier les lettres).
    \item \terme{Déformer} un texte (appliquer des déformations).
    \item \terme{Partager} un fichier avec quelqu'un qui n'a pas la police.
\end{itemize}

\subsubsection{Les cas où la vectorisation n'est pas nécessaire}

\begin{itemize}
    \item \terme{Travailler} sur un texte qui doit rester modifiable.
    \item \terme{Éditer} le texte ultérieurement.
    \item \terme{Conserver} la police d'origine.
\end{itemize}

\subsection{Comment vectoriser un texte ?}

\subsubsection{Par le menu}

\begin{enumerate}
    \item Sélectionner le texte.
    \item Utiliser \texttt{Texte > Vectoriser}.
    \item Le texte devient un groupe de formes.
\end{enumerate}

\subsubsection{Par le raccourci}

\texttt{Ctrl + Maj + O} (ou \texttt{Cmd + Maj + O}).

\subsection{Les conséquences de la vectorisation}

\subsubsection{Le texte n'est plus modifiable}

Une fois vectorisé, le texte ne peut plus être édité comme du texte. Il faut le modifier comme des formes.

\subsubsection{Le texte devient un groupe}

Le texte vectorisé devient un groupe de formes. On peut :

\begin{itemize}
    \item \terme{Dégrouper} le groupe pour modifier chaque lettre.
    \item \terme{Modifier} les points et les courbes.
    \item \terme{Combiner} les formes.
\end{itemize}

\subsubsection{Le texte perd les informations de police}

Le texte vectorisé ne contient plus l'information de la police. On ne peut plus changer la police.

\subsection{Les bonnes pratiques}

\begin{itemize}
    \item \terme{Conserver} une copie du texte non vectorisé (sur un calque séparé).
    \item \terme{Vectoriser} uniquement si nécessaire.
    \item \terme{Vérifier} le résultat après vectorisation.
    \item \terme{Modifier} les formes si nécessaire.
    \item \terme{Incorporer} les polices dans le PDF (alternative à la vectorisation).
    \item \terme{Enregistrer} les deux versions (texte et vectorisé).
\end{itemize}

\subsection*{Conclusion de la section}

La \terme{vectorisation de texte} transforme un texte en formes vectorielles. Elle est nécessaire pour envoyer un fichier à un imprimeur, créer un logo, personnaliser un texte, déformer un texte. Elle a des conséquences : le texte n'est plus modifiable, devient un groupe de formes, perd les informations de police. Il est essentiel de conserver une copie du texte non vectorisé, de vectoriser uniquement si nécessaire, de vérifier le résultat.

Maintenant que nous savons vectoriser un texte, nous allons, dans la section suivante, apprendre à \terme{personnaliser des lettres}.

% ============================================================
\section{Personnaliser des lettres}
% ============================================================

\subsection{Qu'est-ce que la personnalisation de lettres ?}

\begin{quote}
\textbf{La personnalisation de lettres consiste à modifier les lettres d'un texte (ou de créer des lettres uniques) pour leur donner un style original, adapté à un projet, à une marque, à une ambiance.}
\end{quote}

La personnalisation de lettres est au cœur du lettrage. Elle permet de créer des caractères uniques, qui n'existent dans aucune police.

\subsection{Les techniques de personnalisation}

\subsubsection{Modifier les lettres vectorisées}

\begin{enumerate}
    \item Vectoriser le texte.
    \item Dégrouper le groupe.
    \item Modifier chaque lettre (déplacer les points, ajuster les courbes).
    \item Combiner les formes si nécessaire.
\end{enumerate}

\subsubsection{Ajouter des éléments}

\begin{itemize}
    \item \terme{Ajouter} des fioritures.
    \item \terme{Ajouter} des ornements.
    \item \terme{Ajouter} des extensions.
    \item \terme{Ajouter} des effets.
\end{itemize}

\subsubsection{Déformer les lettres}

\begin{itemize}
    \item \terme{Allonger} une lettre.
    \item \terme{Élargir} une lettre.
    \item \terme{Incliner} une lettre.
    \item \terme{Courber} une lettre.
\end{itemize}

\subsubsection{Combiner les lettres}

\begin{itemize}
    \item \terme{Lier} deux lettres (ligature).
    \item \terme{Superposer} deux lettres.
    \item \terme{Entrelacer} deux lettres.
\end{itemize}

\subsubsection{Créer des lettres à partir de zéro}

\begin{enumerate}
    \item Dessiner les lettres avec la Plume.
    \item Ajuster les courbes.
    \item Combiner les formes.
    \item Appliquer les couleurs.
\end{enumerate}

\subsection{Les outils utiles}

\begin{itemize}
    \item \terme{La Plume} : pour dessiner et modifier les lettres.
    \item \terme{Le Pathfinder} : pour combiner des lettres.
    \item \terme{Le Concepteur de forme} : pour fusionner des lettres.
    \item \terme{La Déformation} : pour déformer les lettres.
    \item \terme{L'outil Largeur} : pour varier l'épaisseur des traits.
\end{itemize}

\subsection{Les bonnes pratiques}

\begin{itemize}
    \item \terme{Commencer} par vectoriser le texte.
    \item \terme{Dégrouper} le groupe pour modifier chaque lettre.
    \item \terme{Modifier} avec la Plume et le Pathfinder.
    \item \terme{Combiner} les lettres si nécessaire.
    \item \terme{Vérifier} la lisibilité.
    \item \terme{Conserver} une copie du texte original.
\end{itemize}

\subsection*{Conclusion de la section}

La \terme{personnalisation de lettres} consiste à modifier les lettres d'un texte pour créer des caractères uniques. On utilise la vectorisation, le dégroupage, la modification avec la Plume, le Pathfinder, le Concepteur de forme, la Déformation. Il est essentiel de commencer par vectoriser, de modifier chaque lettre, de combiner, de vérifier la lisibilité.

Maintenant que nous savons personnaliser des lettres, nous allons, dans la section suivante, apprendre à \terme{créer un lettrage simple}.

% ============================================================
\section{Créer un lettrage simple}
% ============================================================

\subsection{Qu'est-ce que le lettrage ?}

\begin{quote}
\textbf{Le lettrage (ou \emph{lettering}) est l'art de dessiner des lettres à la main, de créer des caractères uniques et personnalisés, adaptés à un projet, à une marque, à une ambiance.}
\end{quote}

Le lettrage diffère de la typographie :

\begin{itemize}
    \item \terme{La typographie} utilise des polices existantes.
    \item \terme{Le lettrage} crée des lettres uniques.
    \item \terme{La typographie} est standardisée.
    \item \terme{Le lettrage} est personnalisé.
\end{itemize}

Le lettrage est utilisé pour :

\begin{itemize}
    \item \terme{Logos} typographiques.
    \item \terme{Titres} d'affiches.
    \item \terme{Couvertures} de livres.
    \item \terme{Identités visuelles}.
    \item \terme{Illustrations} avec texte.
\end{itemize}

\subsection{Les étapes de création d'un lettrage}

\subsubsection{Étape 1 : Concevoir}

\begin{enumerate}
    \item \terme{Définir} le mot ou la phrase à lettrer.
    \item \terme{Définir} l'ambiance (élégante, dynamique, créative, etc.).
    \item \terme{Esquisser} plusieurs pistes sur papier.
    \item \terme{Choisir} la piste la plus pertinente.
\end{enumerate}

\subsubsection{Étape 2 : Choisir la structure}

\begin{enumerate}
    \item \terme{Décider} de la forme globale (ligne droite, courbe, cercle).
    \item \terme{Décider} du style (script, sans-serif, serif, décoratif).
    \item \terme{Décider} de l'épaisseur (fine, moyenne, épaisse).
\end{enumerate}

\subsubsection{Étape 3 : Dessiner les lettres de base}

\begin{enumerate}
    \item \terme{Dessiner} les lettres avec la Plume.
    \item \terme{Utiliser} des formes géométriques pour les lettres de base.
    \item \terme{Aligner} les lettres.
\end{enumerate}

\subsubsection{Étape 4 : Personnaliser les lettres}

\begin{enumerate}
    \item \terme{Modifier} les lettres (allonger, élargir, courber).
    \item \terme{Ajouter} des fioritures.
    \item \terme{Combiner} les lettres si nécessaire.
\end{enumerate}

\subsubsection{Étape 5 : Appliquer les couleurs}

\begin{enumerate}
    \item \terme{Appliquer} les couleurs.
    \item \terme{Vérifier} la cohérence.
\end{enumerate}

\subsubsection{Étape 6 : Vérifier}

\begin{enumerate}
    \item \terme{Vérifier} la lisibilité.
    \item \terme{Vérifier} l'esthétique.
    \item \terme{Vérifier} la cohérence.
\end{enumerate}

\subsection{Les techniques de lettrage}

\subsubsection{Le lettrage script}

Le \terme{lettrage script} imite l'écriture manuscrite. Il est fluide, élégant, expressif. Il est utilisé pour les logos, les invitations, les couvertures.

\subsubsection{Le lettrage sans-serif}

Le \terme{lettrage sans-serif} est moderne, géométrique, lisible. Il est utilisé pour les identités visuelles, les affiches, les titres.

\subsubsection{Le lettrage serif}

Le \terme{lettrage serif} est classique, élégant, traditionnel. Il est utilisé pour les logos, les titres, les éditions.

\subsubsection{Le lettrage décoratif}

Le \terme{lettrage décoratif} est original, créatif, expressif. Il est utilisé pour les affiches, les logos, les illustrations.

\subsection{Les bonnes pratiques}

\begin{itemize}
    \item \terme{Commencer} par esquisser sur papier.
    \item \terme{Choisir} une structure claire.
    \item \terme{Dessiner} avec la Plume.
    \item \terme{Personnaliser} les lettres.
    \item \terme{Aligner} et équilibrer.
    \item \terme{Vérifier} la lisibilité.
    \item \terme{Vérifier} la cohérence.
    \item \terme{Conserver} les versions modifiables.
\end{itemize}

\subsection*{Conclusion de la section}

Le \terme{lettrage} est l'art de dessiner des lettres à la main. Il diffère de la typographie (polices existantes) par sa personnalisation. La création suit plusieurs étapes : concevoir, choisir la structure, dessiner les lettres de base, personnaliser, appliquer les couleurs, vérifier. Les styles sont : script, sans-serif, serif, décoratif. Il est essentiel de commencer par esquisser, de dessiner avec la Plume, de personnaliser, de vérifier la lisibilité.

Maintenant que nous savons créer un lettrage simple, nous allons, dans la section suivante, apprendre à \terme{concevoir un titre publicitaire}.

% ============================================================
\section{Concevoir un titre publicitaire}
% ============================================================

\subsection{Qu'est-ce qu'un titre publicitaire ?}

\begin{quote}
\textbf{Un titre publicitaire est un texte court, percutant, mémorable, qui transmet le message principal d'une publicité et attire l'attention.}
\end{quote}

Le titre publicitaire est l'élément le plus visible d'une publicité. Il doit :

\begin{itemize}
    \item \terme{Attirer} l'attention.
    \item \terme{Transmettre} le message principal.
    \item \terme{Créer} une émotion.
    \item \terme{Être} mémorable.
    \item \terme{Inciter} à l'action.
\end{itemize}

\subsection{Les caractéristiques d'un bon titre}

\subsubsection{La concision}

Un bon titre est \terme{court} : 3 à 7 mots en général. Il doit être lu en quelques secondes.

\subsubsection{La clarté}

Un bon titre est \terme{clair} : le message est compris immédiatement.

\subsubsection{L'impact}

Un bon titre est \terme{percutant} : il surprend, il marque, il interpelle.

\subsubsection{La mémorisation}

Un bon titre est \terme{mémorable} : on s'en souvient, on le répète, on le partage.

\subsubsection{La cohérence}

Un bon titre est \terme{cohérent} avec la marque, le produit, le public.

\subsection{La conception du titre}

\subsubsection{Étape 1 : Définir le message}

Quel est le message principal ? Que doit retenir le spectateur ?

\subsubsection{Étape 2 : Choisir la typographie}

Quelle typographie pour le titre ? Adaptée à l'ambiance, lisible, impactante.

\subsubsection{Étape 3 : Composer}

Comment composer le titre ? Hiérarchie, alignement, taille, couleur.

\subsubsection{Étape 4 : Appliquer les effets}

Quels effets appliquer ? Avec parcimonie (ombre, contour, dégradé).

\subsubsection{Étape 5 : Vérifier}

Vérifier la lisibilité, l'impact, la cohérence.

\subsection{Les styles de titres publicitaires}

\subsubsection{Le titre classique}

Le \terme{titre classique} utilise une typographie serif ou sans-serif, un alignement centré, une couleur sobre. Il est utilisé pour les publicités institutionnelles, les marques de luxe.

\subsubsection{Le titre moderne}

Le \terme{titre moderne} utilise une typographie sans-serif, un alignement à gauche, un contraste fort. Il est utilisé pour les publicités tech, les start-up.

\subsubsection{Le titre expressif}

Le \terme{titre expressif} utilise une typographie décorative, des déformations, des effets. Il est utilisé pour les publicités créatives, les événements.

\subsubsection{Le titre minimaliste}

Le \terme{titre minimaliste} utilise une typographie simple, un espacement large, une couleur neutre. Il est utilisé pour les marques de luxe, les publicités épurées.

\subsection{Les bonnes pratiques}

\begin{itemize}
    \item \terme{Un seul message} principal.
    \item \terme{Un titre court} et percutant.
    \item \terme{Une typographie} adaptée.
    \item \terme{Une hiérarchie} claire.
    \item \terme{Des effets} avec parcimonie.
    \item \terme{Une lisibilité} maximale.
    \item \terme{Une cohérence} avec la marque.
    \item \terme{Un impact} fort.
\end{itemize}

\subsection*{Conclusion de la section}

Un \terme{titre publicitaire} est un texte court, percutant, mémorable, qui transmet le message principal. Il doit être concis, clair, impactant, mémorable, cohérent. La conception suit plusieurs étapes : définir le message, choisir la typographie, composer, appliquer les effets, vérifier. Les styles sont : classique, moderne, expressif, minimaliste. Il est essentiel d'avoir un seul message, un titre court, une typographie adaptée, une lisibilité maximale.

Maintenant que nous savons concevoir un titre publicitaire, nous allons, dans la dernière section de ce chapitre, apprendre à \terme{préserver les versions modifiables des textes}.

% ============================================================
\section{Préserver les versions modifiables des textes}
% ============================================================

\subsection{Pourquoi préserver les versions modifiables ?}

La \terme{préservation} des versions modifiables des textes est essentielle pour :

\begin{itemize}
    \item \terme{Modifier} le texte ultérieurement.
    \item \terme{Corriger} des erreurs.
    \item \terme{Adapter} le texte à différents supports.
    \item \terme{Travailler} en équipe.
    \item \terme{Reprendre} un projet après un certain temps.
\end{itemize}

Une fois le texte vectorisé, il n'est plus modifiable. Il est donc essentiel de conserver une version modifiable.

\subsection{Les méthodes de préservation}

\subsubsection{Conserver le texte original sur un calque séparé}

\begin{enumerate}
    \item Créer un calque pour le texte original.
    \item Nommer le calque (par exemple, « Texte - Original »).
    \item Verrouiller le calque.
    \item Masquer le calque.
    \item Créer un calque pour le texte vectorisé.
    \item Vectoriser le texte sur ce calque.
\end{enumerate}

\subsubsection{Enregistrer deux versions du fichier}

\begin{enumerate}
    \item Enregistrer une version avec le texte modifiable (\texttt{fichier-texte.ai}).
    \item Enregistrer une version avec le texte vectorisé (\texttt{fichier-vectorise.ai}).
\end{enumerate}

\subsubsection{Utiliser les objets dynamiques}

Les objets dynamiques permettent de conserver les données d'origine d'un texte, même après transformation.

\subsection{Les bonnes pratiques}

\begin{itemize}
    \item \terme{Conserver} une copie du texte original.
    \item \terme{Nommer} les calques de manière claire.
    \item \terme{Verrouiller} les calques terminés.
    \item \terme{Enregistrer} plusieurs versions.
    \item \terme{Utiliser} des objets dynamiques.
    \item \terme{Communiquer} avec l'équipe.
    \item \terme{Documenter} les modifications.
\end{itemize}

\subsection{Les erreurs courantes}

\subsubsection{Vectoriser sans conserver l'original}

Une fois vectorisé, le texte n'est plus modifiable. Il faut toujours conserver une copie.

\subsubsection{Supprimer le calque original}

Le calque original est la seule version modifiable. Il ne faut pas le supprimer.

\subsubsection{Ne pas nommer les calques}

Des calques mal nommés rendent difficile la reprise du projet.

\subsubsection{Ne pas enregistrer de versions}

Une seule version du fichier rend difficile le retour en arrière.

\subsection*{Conclusion de la section}

La \terme{préservation} des versions modifiables des textes est essentielle pour pouvoir modifier, corriger, adapter, travailler en équipe. On conserve le texte original sur un calque séparé, on enregistre plusieurs versions, on utilise des objets dynamiques. Il est essentiel de conserver une copie du texte original, de nommer les calques, d'enregistrer plusieurs versions.

% ============================================================
\section*{Conclusion du chapitre}
\addcontentsline{toc}{section}{Conclusion du chapitre}
% ============================================================

Nous voici au terme du chapitre 19. Nous avons appris à maîtriser la typographie et le lettrage dans Illustrator :

\begin{itemize}
    \item \terme{Créer et modifier du texte} : outils (Texte, Texte vertical, Texte sur tracé), types (ponctuel, paragraphe, curviligne), vectorisation.
    \item \terme{Propriétés typographiques} : police, taille, style, interlignage, crénage, approche, couleur, alignement, styles.
    \item \terme{Crénage et interlignage} : définitions, méthodes, valeurs recommandées, applications.
    \item \terme{Compositions typographiques} : hiérarchie, alignement, contraste, espace ; types (classique, moderne, expressif, minimaliste).
    \item \terme{Texte sur tracé} : création, déplacement, retournement, options.
    \item \terme{Vectorisation} : définition, cas d'usage, conséquences, bonnes pratiques.
    \item \terme{Personnalisation de lettres} : techniques (modifier, ajouter, déformer, combiner, créer).
    \item \terme{Lettrage} : définition, étapes, techniques (script, sans-serif, serif, décoratif).
    \item \terme{Titre publicitaire} : caractéristiques, conception, styles.
    \item \terme{Préservation des versions modifiables} : méthodes, bonnes pratiques.
\end{itemize}

La typographie est l'un des piliers du design graphique. Dans Illustrator, elle prend une dimension particulière : le texte est vectoriel, modifiable, déformable, personnalisable. On peut créer des compositions typographiques précises, des lettrages uniques, des titres publicitaires impactants. Un designer qui maîtrise la typographie dans Illustrator peut créer des visuels typographiques professionnels ; un designer qui ne la maîtrise pas reste limité à des textes approximatifs.

Ce chapitre t'a donné les bases de la typographie et du lettrage dans Illustrator. Mais la maîtrise vient avec la pratique. N'hésite pas à expérimenter, à créer des compositions, à personnaliser des lettres, à dessiner des lettrages. Chaque projet est une occasion de progresser.

Dans le chapitre suivant, nous aborderons la \terme{création de logos professionnels}. Tu apprendras à comprendre la fonction d'un logo, à analyser un brief, à effectuer des recherches visuelles, à esquisser plusieurs concepts, à construire un logo géométrique, un logo typographique, un symbole vectoriel, à évaluer la lisibilité à différentes tailles, à tester les versions monochromes, à préparer les variantes, à organiser les fichiers sources et les exports. Ces compétences te permettront de créer des logos professionnels.

\vspace{2em}

\begin{center}
\textit{Fin du chapitre 19 --- Typographie et lettrage dans Illustrator.}

\vspace{1em}

\textbf{Prochain chapitre : Chapitre 20 --- Création de logos professionnels.}
\end{center}